\documentclass[11pt, a4paper]{article}
\usepackage[utf8]{inputenc}
\usepackage[margin=0.75in]{geometry}
\usepackage{amsmath, amssymb}
\usepackage{enumitem}
\usepackage{multicol}

\title{\textbf{NEET Biology: Fermentation — High-Yield Difficult Questions}}
\author{Target: NEET UG}
\date{}

\begin{document}

\maketitle

\section*{Questions}

\begin{enumerate}[label=\textbf{\arabic*.}]

\item Which of the following statements accurately describes the fundamental biochemical purpose of fermentation in anaerobic organisms and hypoxic tissues?
\begin{enumerate}[label=(\Alph*)]
    \item To generate additional $\text{ATP}$ molecules via substrate-level phosphorylation beyond glycolysis.
    \item To completely oxidize pyruvic acid into $\text{CO}_2$ and $\text{H}_2\text{O}$ in the absence of $\text{O}_2$.
    \item To reoxidize $\text{NADH} + \text{H}^+$ back to $\text{NAD}^+$ so that glycolysis can continue uninterrupted.
    \item To convert toxic lactic acid into non-toxic ethanol and carbon dioxide.
\end{enumerate}

\item According to NCERT, what approximate percentage of the total free energy contained within a glucose molecule is released during lactic acid or alcoholic fermentation, and how much of it is trapped as high-energy phosphate bonds of $\text{ATP}$?
\begin{enumerate}[label=(\Alph*)]
    \item Less than $7\%$ is released, and all of it is trapped as $\text{ATP}$.
    \item Less than $7\%$ is released, and not all of it is trapped as $\text{ATP}$.
    \item About $38\%$ is released, and $7\%$ is trapped as $\text{ATP}$.
    \item Exactly $2\%$ is released, and $100\%$ of it is trapped as $\text{ATP}$.
\end{enumerate}

\item At what approximate concentration of alcohol do fermenting yeast cells (*Saccharomyces cerevisiae*) poison themselves to death?
\begin{enumerate}[label=(\Alph*)]
    \item $5\%$
    \item $9\%$
    \item $13\%$
    \item $21\%$
\end{enumerate}

\item \textbf{Assertion (A):} In lactic acid fermentation, no carbon dioxide ($\text{CO}_2$) is released, whereas alcoholic fermentation releases two molecules of $\text{CO}_2$ per glucose molecule. \\
\textbf{Reason (R):} Lactic acid is a 3-carbon compound formed by the direct reduction of pyruvic acid, whereas ethanol is a 2-carbon compound formed following the non-oxidative decarboxylation of pyruvic acid.
\begin{enumerate}[label=(\Alph*)]
    \item Both (A) and (R) are true and (R) is the correct explanation of (A).
    \item Both (A) and (R) are true but (R) is not the correct explanation of (A).
    \item (A) is true but (R) is false.
    \item Both (A) and (R) are false.
\end{enumerate}

\item Identify the correct sequence of enzymes and cofactors required to convert pyruvic acid into ethanol in yeast:
\begin{enumerate}[label=(\Alph*)]
    \item Pyruvate dehydrogenase ($\text{FAD}, \text{CoA}$) $\rightarrow$ Alcohol dehydrogenase ($\text{NAD}^+$)
    \item Pyruvate decarboxylase ($\text{TPP}, \text{Mg}^{2+}$) $\rightarrow$ Alcohol dehydrogenase ($\text{Zn}^{2+}, \text{NADH} + \text{H}^+$)
    \item Lactate dehydrogenase ($\text{FMN}$) $\rightarrow$ Pyruvate decarboxylase ($\text{Biotin}$)
    \item Pyruvate kinase ($\text{K}^+$) $\rightarrow$ Acetaldehyde reductase ($\text{NADPH}$)
\end{enumerate}

\item In alcoholic fermentation, which molecule serves as the **terminal electron acceptor** for the electrons carried by $\text{NADH} + \text{H}^+$?
\begin{enumerate}[label=(\Alph*)]
    \item Pyruvic acid
    \item Acetaldehyde
    \item Ethanol
    \item Molecular oxygen ($\text{O}_2$)
\end{enumerate}

\item In homolactic fermentation (such as in *Lactobacillus* or strenuously exercising human skeletal muscle), which molecule serves as the **final hydrogen/electron acceptor**?
\begin{enumerate}[label=(\Alph*)]
    \item Glyceraldehyde-3-phosphate
    \item Acetaldehyde
    \item Pyruvic acid
    \item Lactic acid
\end{enumerate}

\item If a glucose molecule is radiolabeled with $^{14}\text{C}$ exclusively at Carbon-2 ($\text{C}_2$) and subjected to complete alcoholic fermentation by yeast, where will the radioactive label be located at the end of the process?
\begin{enumerate}[label=(\Alph*)]
    \item Exclusively in the $\text{CO}_2$ evolved
    \item Exclusively in the methyl carbon ($-\text{CH}_3$) of ethanol
    \item Exclusively in the carbinol carbon ($-\text{CH}_2\text{OH}$) of ethanol
    \item Equally distributed between $-\text{CH}_2\text{OH}$ of ethanol and $\text{CO}_2$
\end{enumerate}

\item Consider the following statements comparing Lactic Acid Fermentation (LAF) and Alcoholic Fermentation (AF):
\begin{enumerate}[label=(\roman*)]
    \item Both processes yield a net gain of $2\text{ ATP}$ molecules per molecule of glucose oxidized.
    \item Both processes involve a net consumption of $2\text{ NADH} + \text{H}^+$ during the reduction of pyruvate or its derivative.
    \item Pyruvate decarboxylase is present in human skeletal muscle cells but remains inactive under aerobic conditions.
    \item Both processes are hazardous because either acid or alcohol is produced.
\end{enumerate}
Which of the above statements are correct?
\begin{enumerate}[label=(\Alph*)]
    \item (i), (ii), and (iv)
    \item (i) and (iii) only
    \item (ii), (iii), and (iv)
    \item (i), (ii), (iii), and (iv)
\end{enumerate}

\item Why can humans never produce ethanol in their skeletal muscles during extreme oxygen deprivation, even though all the necessary glycolytic enzymes are present?
\begin{enumerate}[label=(\Alph*)]
    \item Human muscle cells lack the enzyme alcohol dehydrogenase.
    \item Human muscle cells lack the enzyme pyruvate decarboxylase.
    \item Acetaldehyde is immediately converted into acetyl-CoA in human cytosol.
    \item Human lactate dehydrogenase has a higher affinity for acetaldehyde than for pyruvate.
\end{enumerate}

\item Match Column-I (Organism / Tissue) with Column-II (Primary Fermentation Product & Key Enzyme):
\begin{center}
\begin{tabular}{|l|l|}
\hline
\textbf{Column-I} & \textbf{Column-II} \\
\hline
P. *Lactobacillus bulgaricus* & 1. Ethanol + $\text{CO}_2$; Pyruvate decarboxylase \\
Q. *Saccharomyces cerevisiae* & 2. Lactic acid; Lactate dehydrogenase \\
R. *Propionibacterium shermanii* & 3. Propionic acid + Acetic acid + $\text{CO}_2$ \\
S. Waterlogged roots of higher plants & 4. Ethanol + $\text{CO}_2$ (transiently) \\
\hline
\end{tabular}
\end{center}
\begin{enumerate}[label=(\Alph*)]
    \item P-2, Q-1, R-3, S-4
    \item P-1, Q-2, R-3, S-4
    \item P-2, Q-3, R-1, S-4
    \item P-4, Q-1, R-2, S-3
\end{enumerate}

\item During strenuous physical exercise, lactic acid accumulates in human skeletal muscles. What happens to the majority of this lactate once oxygen becomes available again during recovery (the Cori Cycle)?
\begin{enumerate}[label=(\Alph*)]
    \item It is excreted unchanged in urine and sweat.
    \item It is decarboxylated directly into ethanol in the kidneys.
    \item It is transported via blood to the liver, oxidized back to pyruvate, and converted into glucose via gluconeogenesis.
    \item It enters the mitochondrial matrix of RBCs to run the Krebs cycle.
\end{enumerate}

\item How many total molecules of $\text{NADH} + \text{H}^+$ are **net** produced at the end of complete fermentation of $3\text{ molecules}$ of glucose into lactic acid?
\begin{enumerate}[label=(\Alph*)]
    \item 0
    \item 3
    \item 6
    \item 12
\end{enumerate}

\item \textbf{Statement I:} In fermentation, there is a net gain of only $2\text{ ATP}$ molecules for each molecule of glucose degraded to pyruvic acid. \\
\textbf{Statement II:} NADH is oxidized to $\text{NAD}^+$ rather slowly in fermentation; however, the reaction is very vigorous in case of aerobic respiration.
\begin{enumerate}[label=(\Alph*)]
    \item Both Statement I and Statement II are correct.
    \item Both Statement I and Statement II are incorrect.
    \item Statement I is correct, but Statement II is incorrect.
    \item Statement I is incorrect, but Statement II is correct.
\end{enumerate}

\item What is the fate of the pyruvate carbons when $1\text{ molecule}$ of glucose labeled at $\text{C}_1$ and $\text{C}_6$ undergoes lactic acid fermentation?
\begin{enumerate}[label=(\Alph*)]
    \item Both labeled carbons are lost as $\text{CO}_2$.
    \item Both labeled carbons appear in the carboxyl group ($-\text{COOH}$) of lactic acid.
    \item Both labeled carbons appear in the methyl group ($-\text{CH}_3$) of lactic acid.
    \item One appears in the $-\text{CHOH}-$ group and one in the $-\text{COOH}$ group of lactic acid.
\end{enumerate}

\item Which of the following best explains why alcoholic beverages with an alcohol content greater than $13\%$ (such as brandy, whisky, or rum) cannot be produced solely by natural yeast fermentation?
\begin{enumerate}[label=(\Alph*)]
    \item Glucose becomes completely insoluble in water above $13\%$ ethanol.
    \item Ethanol denatures membrane lipids and key enzymes of yeast above $\sim 13\%$, killing the cells; higher concentrations require distillation.
    \item Pyruvate decarboxylase is allosterically inhibited by $\text{CO}_2$ when ethanol reaches $13\%$.
    \item Above $13\%$ ethanol, yeast switches to aerobic respiration using dissolved $\text{CO}_2$.
\end{enumerate}

\item In the conversion of pyruvic acid to acetaldehyde catalyzed by pyruvate decarboxylase, which vitamin-derived coenzyme is strictly essential?
\begin{enumerate}[label=(\Alph*)]
    \item Riboflavin ($\text{FAD}$)
    \item Niacin ($\text{NAD}^+$)
    \item Thiamine pyrophosphate ($\text{TPP}$, from Vitamin $\text{B}_1$)
    \item Pyridoxal phosphate ($\text{PLP}$, from Vitamin $\text{B}_6$)
\end{enumerate}

\item Which of the following enzymes involved in alcoholic fermentation is a **zinc-containing metalloenzyme**?
\begin{enumerate}[label=(\Alph*)]
    \item Hexokinase
    \item Enolase
    \item Alcohol dehydrogenase
    \item Pyruvate decarboxylase
\end{enumerate}

\item Compare the net $\text{ATP}$ yield per molecule of glucose in a facultative anaerobe (like *Saccharomyces cerevisiae*) shifting from strict anaerobic alcoholic fermentation to complete aerobic respiration (via the malate-aspartate shuttle). By what factor does the net $\text{ATP}$ yield increase?
\begin{enumerate}[label=(\Alph*)]
    \item $2\text{-fold}$
    \item $15\text{-fold}$
    \item $18\text{-fold}$
    \item $19\text{-fold}$
\end{enumerate}

\item Why does the Pasteur effect occur when a culture of facultatively anaerobic yeast growing under anaerobic conditions is suddenly aerated?
\begin{enumerate}[label=(\Alph*)]
    \item Oxygen directly binds to and inhibits alcohol dehydrogenase.
    \item Aerobic respiration produces vastly more $\text{ATP}$ per glucose ($36\text{--}38\text{ ATP}$ vs. $2\text{ ATP}$), and elevated $\text{ATP}$ and citrate allosterically inhibit phosphofructokinase-1 ($\text{PFK-1}$), slowing down glucose consumption.
    \item Aeration lowers the pH of the medium, denaturing hexokinase.
    \item Oxygen oxidizes ethanol back into glucose via reverse glycolysis.
\end{enumerate}

\item Which of the following reactions represents the exact stoichiometry of alcoholic fermentation of one molecule of glucose?
\begin{enumerate}[label=(\Alph*)]
    \item $\text{C}_6\text{H}_{12}\text{O}_6 + 2\text{ADP} + 2\text{P}_i \rightarrow 2\text{C}_2\text{H}_5\text{OH} + 2\text{CO}_2 + 2\text{ATP} + 2\text{H}_2\text{O}$
    \item $\text{C}_6\text{H}_{12}\text{O}_6 + 2\text{ADP} + 2\text{P}_i + 2\text{NAD}^+ \rightarrow 2\text{C}_2\text{H}_5\text{OH} + 2\text{CO}_2 + 2\text{ATP} + 2\text{NADH} + 2\text{H}^+$
    \item $\text{C}_6\text{H}_{12}\text{O}_6 + 2\text{ATP} \rightarrow 2\text{CH}_3\text{CH(OH)COOH} + 2\text{ADP} + 2\text{P}_i$
    \item $\text{C}_6\text{H}_{12}\text{O}_6 + 6\text{O}_2 \rightarrow 2\text{C}_2\text{H}_5\text{OH} + 4\text{CO}_2 + 2\text{H}_2\text{O}$
\end{enumerate}

\item Mature human erythrocytes (RBCs) lack mitochondria and therefore rely exclusively on glycolysis followed by lactic acid fermentation for energy. How many $\text{ATP}$ molecules are generated per glucose molecule metabolized in an RBC (ignoring the Rapoport-Luebering shunt)?
\begin{enumerate}[label=(\Alph*)]
    \item $0\text{ ATP}$
    \item $2\text{ ATP}$
    \item $4\text{ ATP}$
    \item $8\text{ ATP}$
\end{enumerate}

\item What is the primary difference between **fermentation** and **anaerobic respiration** in strict biochemical terms?
\begin{enumerate}[label=(\Alph*)]
    \item Fermentation uses molecular oxygen, whereas anaerobic respiration does not.
    \item Fermentation uses an endogenous organic molecule (like pyruvate or acetaldehyde) as the final electron acceptor without an electron transport chain, whereas anaerobic respiration uses an exogenous inorganic electron acceptor (like $\text{NO}_3^-$, $\text{SO}_4^{2-}$, or $\text{CO}_2$) via a membrane-bound electron transport chain.
    \item Anaerobic respiration occurs only in eukaryotes, while fermentation occurs only in prokaryotes.
    \item Fermentation yields $38\text{ ATP}$, whereas anaerobic respiration yields $2\text{ ATP}$.
\end{enumerate}

\item If a yeast culture is supplied with $1\text{ mole}$ of **sucrose** under strictly anaerobic conditions, how many moles of $\text{CO}_2$ and ethanol are produced upon complete fermentation?
\begin{enumerate}[label=(\Alph*)]
    \item $2\text{ moles}$ of $\text{CO}_2$ and $2\text{ moles}$ of ethanol
    \item $4\text{ moles}$ of $\text{CO}_2$ and $4\text{ moles}$ of ethanol
    \item $6\text{ moles}$ of $\text{CO}_2$ and $2\text{ moles}$ of ethanol
    \item $4\text{ moles}$ of $\text{CO}_2$ and $2\text{ moles}$ of ethanol
\end{enumerate}

\item \textbf{Assertion (A):} In lactic acid fermentation, the $\text{NADH} + \text{H}^+$ generated during the oxidation of glyceraldehyde-3-phosphate is reoxidized to $\text{NAD}^+$ when pyruvic acid is reduced to lactic acid. \\
\textbf{Reason (R):} The reduction of pyruvic acid to lactic acid is an endergonic reaction that requires the hydrolysis of $1\text{ ATP}$ per pyruvate molecule.
\begin{enumerate}[label=(\Alph*)]
    \item Both (A) and (R) are true and (R) is the correct explanation of (A).
    \item Both (A) and (R) are true but (R) is not the correct explanation of (A).
    \item (A) is true but (R) is false.
    \item Both (A) and (R) are false.
\end{enumerate}

\item Why does dough used for making *idli*, *dosa*, or bread puff up during fermentation?
\begin{enumerate}[label=(\Alph*)]
    \item Due to rapid evaporation of ethanol produced by *Lactobacillus*.
    \item Due to the production and trapping of $\text{CO}_2$ gas released during fermentation.
    \item Due to the breakdown of gluten proteins into $\text{O}_2$ gas.
    \item Due to osmotic swelling of starch granules in the presence of lactic acid.
\end{enumerate}

\item Which of the following organisms is an **obligate anaerobe** that is poisoned by molecular oxygen and relies exclusively on fermentation or anaerobic pathways?
\begin{enumerate}[label=(\Alph*)]
    \item *Saccharomyces cerevisiae*
    \item *Escherichia coli*
    \item *Clostridium botulinum*
    \item *Pseudomonas aeruginosa*
\end{enumerate}

\item During alcoholic fermentation, the decarboxylation of pyruvic acid to acetaldehyde is an **irreversible** reaction under physiological conditions primarily because:
\begin{enumerate}[label=(\Alph*)]
    \item Alcohol dehydrogenase immediately phosphorylates acetaldehyde.
    \item The gaseous product $\text{CO}_2$ escapes from the aqueous cellular environment, driving the reaction forward by Le Chatelier's principle.
    \item Pyruvate decarboxylase is inhibited by $\text{NAD}^+$.
    \item The reaction consumes $2\text{ ATP}$ molecules.
\end{enumerate}

\item In Swiss cheese production, the characteristic large holes are formed due to the production of a large amount of $\text{CO}_2$ by which bacterium?
\begin{enumerate}[label=(\Alph*)]
    \item *Lactobacillus acidophilus*
    \item *Propionibacterium shermanii*
    \item *Streptococcus thermophilus*
    \item *Acetobacter aceti*
\end{enumerate}

\item How many oxidation-reduction (redox) steps occur in the pathway starting from **glucose** and ending in **ethanol**?
\begin{enumerate}[label=(\Alph*)]
    \item One (Glyceraldehyde-3-phosphate dehydrogenase only)
    \item Two (Glyceraldehyde-3-phosphate dehydrogenase and Alcohol dehydrogenase)
    \item Three (GAPDH, Pyruvate decarboxylase, and Alcohol dehydrogenase)
    \item Zero
\end{enumerate}

\item What is the oxidation number of the carbonyl carbon ($\text{C}_2$) of pyruvic acid before and after its reduction to lactic acid by lactate dehydrogenase?
\begin{enumerate}[label=(\Alph*)]
    \item Changes from $+2$ (ketone) to $0$ (secondary alcohol)
    \item Changes from $+3$ (carboxylic acid) to $+1$ (aldehyde)
    \item Changes from $0$ to $-2$
    \item Remains unchanged at $+2$
\end{enumerate}

\item Which of the following correctly compares the enzyme **Lactate Dehydrogenase (LDH)** isoenzymes $\text{LDH}_5$ (M$_4$, predominant in skeletal muscle) and $\text{LDH}_1$ (H$_4$, predominant in heart muscle)?
\begin{enumerate}[label=(\Alph*)]
    \item $\text{LDH}_5$ is inhibited by high pyruvate levels, whereas $\text{LDH}_1$ rapidly converts pyruvate to lactate.
    \item $\text{LDH}_5$ favors the reduction of pyruvate to lactate under anaerobic bursts, whereas $\text{LDH}_1$ is inhibited by high pyruvate and favors the oxidation of lactate back to pyruvate for aerobic metabolism in the heart.
    \item Both isoenzymes are located exclusively inside the inner mitochondrial membrane.
    \item $\text{LDH}_1$ produces ethanol, while $\text{LDH}_5$ produces lactic acid.
\end{enumerate}

\item Suppose an inhibitor specifically blocks **Alcohol Dehydrogenase** in a yeast cell culture under strictly anaerobic conditions. Which intermediate will immediately accumulate inside the yeast cells before glycolysis halts?
\begin{enumerate}[label=(\Alph*)]
    \item Lactic acid
    \item Acetaldehyde and $\text{NADH} + \text{H}^+$
    \item Acetyl-CoA and $\text{FADH}_2$
    \item Citric acid
\end{enumerate}

\item If a student adds sodium arsenate ($\text{AsO}_4^{3-}$) to a culture of fermenting yeast (which uncouples $\text{ATP}$ synthesis at the phosphoglycerate kinase step without stopping electron flow), what will be the **net $\text{ATP}$ yield** per molecule of glucose fermented?
\begin{enumerate}[label=(\Alph*)]
    \item $2\text{ ATP}$
    \item $1\text{ ATP}$
    \item $0\text{ ATP}$
    \item $-2\text{ ATP}$
\end{enumerate}

\item \textbf{Statement I:} In both lactic acid and alcoholic fermentation, the reducing agent used in the final step is $\text{NADH} + \text{H}^+$ generated during glycolysis. \\
\textbf{Statement II:} In aerobic respiration, the $\text{NADH} + \text{H}^+$ generated during glycolysis is oxidized via the mitochondrial electron transport chain rather than by reducing pyruvate.
\begin{enumerate}[label=(\Alph*)]
    \item Both Statement I and Statement II are correct.
    \item Both Statement I and Statement II are incorrect.
    \item Statement I is correct, but Statement II is incorrect.
    \item Statement I is incorrect, but Statement II is correct.
\end{enumerate}

\item During prolonged waterlogging of agricultural fields, crop plants often die even when soil nutrients are abundant. Which combination of factors best explains this phenomenon at the cellular level?
\begin{enumerate}[label=(\Alph*)]
    \item Root cells switch to alcoholic fermentation, leading to severe $\text{ATP}$ starvation ($2\text{ ATP}$ vs. $38\text{ ATP}$ per glucose) and cytoplasmic toxicity from accumulating ethanol and acetaldehyde.
    \item Root cells accumulate lactic acid, which permanently denatures RuBisCO in the leaves.
    \item Excess water inhibits invertase, preventing sucrose transport from leaves to roots.
    \item Waterlogging stimulates the Krebs cycle excessively, depleting all cellular glucose within minutes.
\end{enumerate}

\item How many molecules of $\text{H}_2\text{O}$ are net produced when $1\text{ molecule}$ of glucose undergoes complete alcoholic fermentation to ethanol and $\text{CO}_2$?
\begin{enumerate}[label=(\Alph*)]
    \item 0
    \item 1
    \item 2
    \item 4
\end{enumerate}

\item Which of the following industrial processes relies on **heterolactic** or mixed-acid fermentation rather than pure homolactic fermentation?
\begin{enumerate}[label=(\Alph*)]
    \item Curdling of milk by *Lactobacillus bulgaricus*
    \item Production of sauerkraut and kimchi involving *Leuconostoc mesenteroides* (producing lactic acid, ethanol, and $\text{CO}_2$)
    \item Vigorous sprint metabolism in human fast-twitch muscle fibers
    \item Glycolysis in mature human red blood cells
\end{enumerate}

\item In the Crabtree effect observed in *Saccharomyces cerevisiae* (baker's yeast), high external concentrations of glucose trigger alcoholic fermentation **even in the presence of abundant oxygen**. What is the primary metabolic reason for this?
\begin{enumerate}[label=(\Alph*)]
    \item High glucose represses the synthesis of mitochondrial respiratory enzymes and saturates the pyruvate dehydrogenase complex, causing overflow of pyruvate into pyruvate decarboxylase.
    \item Oxygen inhibits hexokinase at high glucose concentrations.
    \item Yeast lacks a functional electron transport chain when glucose exceeds $1\text{ mM}$.
    \item Pyruvate kinase converts glucose directly into ethanol without forming pyruvate.
\end{enumerate}

\item Match Column-I (Reaction in Fermentation) with Column-II (Enzyme Class according to IUBMB):
\begin{center}
\begin{tabular}{|l|l|}
\hline
\textbf{Column-I} & \textbf{Column-II} \\
\hline
P. Pyruvic acid $\rightarrow$ Acetaldehyde + $\text{CO}_2$ & 1. Oxidoreductase \\
Q. Acetaldehyde + $\text{NADH} + \text{H}^+$ $\rightarrow$ Ethanol + $\text{NAD}^+$ & 2. Lyase \\
R. Pyruvic acid + $\text{NADH} + \text{H}^+$ $\rightarrow$ Lactic acid + $\text{NAD}^+$ & 3. Transferase \\
\hline
\end{tabular}
\end{center}
\begin{enumerate}[label=(\Alph*)]
    \item P-2, Q-1, R-1
    \item P-1, Q-2, R-3
    \item P-3, Q-1, R-2
    \item P-2, Q-3, R-1
\end{enumerate}

\item Why does human skeletal muscle experience fatigue and localized burning/cramping sensations during intense anaerobic sprinting?
\begin{enumerate}[label=(\Alph*)]
    \item Accumulation of ethanol in the sarcoplasmic reticulum.
    \item Accumulation of lactic acid (dissociating into lactate and $\text{H}^+$ ions), which lowers intracellular pH, inhibits key glycolytic enzymes like $\text{PFK-1}$, and interferes with $\text{Ca}^{2+}$ binding to troponin.
    \item Complete depletion of $\text{CO}_2$ inside muscle fibers.
    \item Excess synthesis of $\text{ATP}$ causing myosin heads to lock onto actin.
\end{enumerate}

\item If $4\text{ molecules}$ of fructose-1,6-bisphosphate are directly fed into an anaerobic cell-free yeast extract capable of glycolysis and alcoholic fermentation, what will be the net yield of $\text{ATP}$, $\text{CO}_2$, and ethanol, respectively?
\begin{enumerate}[label=(\Alph*)]
    \item $8\text{ ATP}$, $8\text{ CO}_2$, $8\text{ Ethanol}$
    \item $16\text{ ATP}$, $8\text{ CO}_2$, $8\text{ Ethanol}$
    \item $8\text{ ATP}$, $4\text{ CO}_2$, $4\text{ Ethanol}$
    \item $16\text{ ATP}$, $16\text{ CO}_2$, $16\text{ Ethanol}$
\end{enumerate}

\item Which of the following statements about fermentation is **incorrect**?
\begin{enumerate}[label=(\Alph*)]
    \item It takes place under anaerobic conditions in many prokaryotes and unicellular eukaryotes.
    \item Incomplete oxidation of glucose is achieved under anaerobic conditions by sets of reactions where pyruvic acid is converted to $\text{CO}_2$ and ethanol or lactic acid.
    \item The reducing agent is $\text{NADH} + \text{H}^+$ which is reoxidized to $\text{NAD}^+$ in both lactic acid and alcohol fermentation.
    \item Oxidation of $\text{NADH}$ to $\text{NAD}^+$ is much faster in fermentation than in aerobic respiration.
\end{enumerate}

\item Consider a facultative anaerobe metabolizing glucose. Under aerobic conditions, $1\text{ molecule}$ of glucose yields $38\text{ ATP}$. Under anaerobic fermentation, $1\text{ molecule}$ of glucose yields $2\text{ ATP}$. To maintain the **exact same rate of $\text{ATP}$ production** ($38\text{ ATP/sec}$) after switching from aerobic to anaerobic conditions, by what factor must the cell increase its rate of glucose consumption?
\begin{enumerate}[label=(\Alph*)]
    \item $2\text{ times}$
    \item $10\text{ times}$
    \item $19\text{ times}$
    \item $38\text{ times}$
\end{enumerate}

\item In *Clostridium acetobutylicum* (used in industrial ABE fermentation), pyruvate is fermented into which commercial solvents?
\begin{enumerate}[label=(\Alph*)]
    \item Acetone, Butanol, and Ethanol
    \item Acetic acid, Butyric acid, and Ethane
    \item Alanine, Biotin, and Ergosterol
    \item Acetaldehyde, Benzene, and Ether
\end{enumerate}

\end{enumerate}

\newpage

\section*{Answer Key}

\begin{multicols}{5}
\begin{enumerate}[label=\textbf{\arabic*.}]
    \item (C)
    \item (B)
    \item (C)
    \item (A)
    \item (B)
    \item (B)
    \item (C)
    \item (C)
    \item (A)
    \item (B)
    \item (A)
    \item (C)
    \item (A)
    \item (A)
    \item (C)
    \item (B)
    \item (C)
    \item (C)
    \item (D)
    \item (B)
    \item (A)
    \item (B)
    \item (B)
    \item (B)
    \item (C)
    \item (B)
    \item (C)
    \item (B)
    \item (B)
    \item (B)
    \item (A)
    \item (B)
    \item (B)
    \item (C)
    \item (A)
    \item (A)
    \item (C)
    \item (B)
    \item (A)
    \item (A)
    \item (B)
    \item (B)
    \item (D)
    \item (C)
    \item (A)
\end{enumerate}
\end{multicols}

\section*{Detailed Solutions}

\begin{enumerate}[label=\textbf{\arabic*.}]

\item \textbf{(C)} Fermentation produces zero $\text{ATP}$ of its own. Its sole metabolic purpose is to regenerate $\text{NAD}^+$ by oxidizing $\text{NADH} + \text{H}^+$, preventing the depletion of the limited cytosolic $\text{NAD}^+$ pool needed by glyceraldehyde-3-phosphate dehydrogenase in glycolysis.

\item \textbf{(B)} Verbatim from NCERT Class 11 Biology (Respiration in Plants): \textit{``Less than seven per cent of the energy in glucose is released in both lactic acid and alcohol fermentation, and not all of it is trapped as high energy bonds of ATP.''}

\item \textbf{(C)} According to NCERT, yeasts poison themselves to death when the concentration of alcohol reaches about **$13\%$**.

\item \textbf{(A)} Pyruvic acid ($3\text{C}$) is directly reduced to lactic acid ($3\text{C}$) by lactate dehydrogenase without carbon loss. In alcoholic fermentation, pyruvate decarboxylase first cleaves off $\text{CO}_2$ ($1\text{C}$) to yield acetaldehyde ($2\text{C}$), which is then reduced to ethanol ($2\text{C}$).

\item \textbf{(B)} Step 1: Pyruvate decarboxylase requires thiamine pyrophosphate ($\text{TPP}$) and $\text{Mg}^{2+}$ to convert pyruvate into acetaldehyde + $\text{CO}_2$. Step 2: Alcohol dehydrogenase requires $\text{Zn}^{2+}$ and $\text{NADH} + \text{H}^+$ to reduce acetaldehyde to ethanol.

\item \textbf{(B)} In alcoholic fermentation, $\text{NADH} + \text{H}^+$ donates its electrons directly to **acetaldehyde** (reducing it to ethanol), making acetaldehyde the terminal electron acceptor.

\item \textbf{(C)} In lactic acid fermentation, **pyruvic acid** directly accepts electrons and protons from $\text{NADH} + \text{H}^+$ to form lactic acid, making pyruvic acid the terminal electron acceptor.

\item \textbf{(C)} During glycolysis, $\text{C}_2$ and $\text{C}_5$ of glucose become the carbonyl carbon ($\text{C}_2$, $-\text{C}=\text{O}$) of the two pyruvate molecules. When pyruvate loses $\text{C}_1$ as $\text{CO}_2$, the carbonyl carbon ($\text{C}_2$) becomes the $-\text{CHO}$ carbon of acetaldehyde, which is reduced to the carbinol carbon ($-\text{CH}_2\text{OH}$) of ethanol.

\item \textbf{(A)} Statements (i), (ii), and (iv) are correct. Statement (iii) is incorrect because humans completely lack the gene/enzyme for pyruvate decarboxylase.

\item \textbf{(B)} Human tissues possess alcohol dehydrogenase (used in the liver to detoxify ingested ethanol), but completely lack **pyruvate decarboxylase**, preventing the conversion of pyruvate into acetaldehyde.

\item \textbf{(A)} *Lactobacillus* produces lactic acid via $\text{LDH}$ (P-2). *Saccharomyces cerevisiae* produces ethanol + $\text{CO}_2$ via pyruvate decarboxylase and $\text{ADH}$ (Q-1). *Propionibacterium shermanii* produces propionic acid, acetic acid, and $\text{CO}_2$ (R-3). Waterlogged plant roots transiently carry out alcoholic fermentation (S-4).

\item \textbf{(C)} Under the **Cori Cycle**, lactate diffuses out of muscles into the bloodstream, travels to the liver, is oxidized back to pyruvate by $\text{LDH}$, and is converted into glucose via gluconeogenesis (consuming $6\text{ ATP}$ in the liver).

\item \textbf{(A)} Glycolysis produces $2\text{ NADH} + \text{H}^+$ per glucose ($6$ for $3\text{ glucose}$), but fermentation consumes the exact same amount ($6\text{ NADH} + \text{H}^+$) to reduce pyruvate to lactate. Net $\text{NADH} + \text{H}^+$ yield = **$0$**.

\item \textbf{(A)} Both statements are direct facts from NCERT: net gain is $2\text{ ATP}$ per glucose degraded to pyruvic acid, and $\text{NADH}$ is oxidized to $\text{NAD}^+$ rather slowly in fermentation compared to the vigorous oxidation in aerobic respiration.

\item \textbf{(C)} In glycolysis, $\text{C}_1$ and $\text{C}_6$ of glucose become the methyl carbons ($\text{C}_3$, $-\text{CH}_3$) of the two pyruvate molecules. Reduction of the keto group at $\text{C}_2$ to $-\text{CHOH}-$ leaves $\text{C}_3$ unchanged as the methyl group ($-\text{CH}_3$) of lactic acid.

\item \textbf{(B)} Ethanol is amphipathic and disrupts the integrity of yeast plasma membranes and denatures metabolic enzymes once it reaches $\sim 13\%$. Higher concentrations in distilled spirits (whisky, rum, brandy) are achieved by physical distillation.

\item \textbf{(C)} Pyruvate decarboxylase requires **Thiamine pyrophosphate ($\text{TPP}$)**, the active coenzyme form of Vitamin $\text{B}_1$ (thiamine), to stabilize the carbanion intermediate during decarboxylation.

\item \textbf{(C)} **Alcohol dehydrogenase ($\text{ADH}$)** contains a catalytic $\text{Zn}^{2+}$ ion at its active site that coordinates the carbonyl oxygen of acetaldehyde, polarizing the $\text{C}=\text{O}$ bond to facilitate hydride transfer from $\text{NADH}$.

\item \textbf{(D)} Anaerobic fermentation yields $2\text{ ATP}$ per glucose. Complete aerobic respiration (with malate-aspartate shuttle) yields $38\text{ ATP}$ per glucose. Increase factor = $38 / 2 = \mathbf{19\text{-fold}}$.

\item \textbf{(B)} Upon aeration, oxidative phosphorylation rapidly generates large amounts of $\text{ATP}$ and citrate. Both $\text{ATP}$ and citrate act as potent allosteric inhibitors of phosphofructokinase-1 ($\text{PFK-1}$), throttling down glycolysis (the Pasteur effect).

\item \textbf{(A)} In glycolysis, $2\text{ ADP} + 2\text{ P}_i \rightarrow 2\text{ ATP} + 2\text{ H}_2\text{O}$ (condensation reaction releases $1\text{ H}_2\text{O}$ per $\text{ATP}$ formed). Meanwhile, $\text{NADH} + \text{H}^+$ produced in glycolysis is completely recycled back to $\text{NAD}^+$, so $\text{NAD}^+/\text{NADH}$ cancels out of the net equation.

\item \textbf{(B)} Mature RBCs lack mitochondria and rely solely on glycolysis + lactic acid fermentation, yielding a net of **$2\text{ ATP}$** per glucose molecule.

\item \textbf{(B)} Fermentation occurs entirely in the cytosol without an electron transport chain, using an endogenous organic molecule (pyruvate/acetaldehyde) as the final electron acceptor. Anaerobic respiration uses an ETC with an external inorganic terminal acceptor ($\text{NO}_3^-$, $\text{SO}_4^{2-}$, $\text{CO}_2$, $\text{Fe}^{3+}$).

\item \textbf{(B)} One mole of sucrose ($12\text{C}$) is hydrolyzed into $1\text{ mole}$ of glucose + $1\text{ mole}$ of fructose ($2\text{ hexoses}$). Each hexose yields $2\text{ moles}$ of $\text{CO}_2$ and $2\text{ moles}$ of ethanol. Thus, $1\text{ mole}$ of sucrose yields $4\text{ moles}$ of $\text{CO}_2$ and $4\text{ moles}$ of ethanol.

\item \textbf{(C)} Assertion (A) is true. Reason (R) is false because the reduction of pyruvate to lactate by $\text{LDH}$ is thermodynamically favorable (exergonic, $\Delta G^{\circ\prime} = -25.1\text{ kJ/mol}$) and does not consume $\text{ATP}$.

\item \textbf{(B)} As noted in NCERT (Microbes in Human Welfare), the puffed-up appearance of dough used for making *dosa*, *idli*, and bread is due to the production of **$\text{CO}_2$ gas** during fermentation.

\item \textbf{(C)} *Clostridium botulinum* is a strict/obligate anaerobe that lacks superoxide dismutase and catalase, making oxygen lethal to it. *S. cerevisiae* and *E. coli* are facultative anaerobes, and *P. aeruginosa* is primarily aerobic (or denitrifying anaerobe).

\item \textbf{(B)} Pyruvate decarboxylase catalyzes a non-oxidative decarboxylation where $\text{CO}_2$ gas escapes from solution, making the reverse carboxylation thermodynamically impossible under cellular conditions.

\item \textbf{(B)} According to NCERT, the large holes in Swiss cheese are due to the production of a large amount of $\text{CO}_2$ by a bacterium named ***Propionibacterium shermanii***.

\item \textbf{(B)} Two redox steps occur: (1) Oxidation of $\text{PGAL}$ to 1,3-bisphosphoglycerate (reducing $\text{NAD}^+$ to $\text{NADH} + \text{H}^+$) in glycolysis, and (2) Reduction of acetaldehyde to ethanol (oxidizing $\text{NADH} + \text{H}^+$ to $\text{NAD}^+$) in fermentation.

\item \textbf{(A)} In pyruvic acid ($\text{CH}_3-\mathbf{C}(=\text{O})-\text{COOH}$), the carbonyl carbon ($\text{C}_2$) is bonded to two carbons and double-bonded to one oxygen (oxidation state = **$+2$**). In lactic acid ($\text{CH}_3-\mathbf{CH}(\text{OH})-\text{COOH}$), $\text{C}_2$ is bonded to two carbons, one $-\text{H}$, and one $-\text{OH}$ (oxidation state = **$0$**).

\item \textbf{(B)} $\text{LDH}_5$ ($\text{M}_4$) in skeletal muscle is adapted for rapid anaerobic reduction of pyruvate to lactate during sprinting. Conversely, the heart is strictly aerobic; $\text{LDH}_1$ ($\text{H}_4$) is inhibited by high pyruvate and preferentially oxidizes lactate into pyruvate to fuel the Krebs cycle.

\item \textbf{(B)} Blocking alcohol dehydrogenase stops the conversion of acetaldehyde to ethanol and prevents the regeneration of $\text{NAD}^+$. Consequently, **acetaldehyde** and **$\text{NADH} + \text{H}^+$** accumulate until $\text{NAD}^+$ exhaustion halts glycolysis.

\item \textbf{(C)} Normal glycolysis invests $2\text{ ATP}$ and generates $4\text{ ATP}$ ($2$ at phosphoglycerate kinase + $2$ at pyruvate kinase). Arsenate abolishes the $2\text{ ATP}$ formed at phosphoglycerate kinase, leaving only $2\text{ ATP}$ formed at pyruvate kinase. Net $\text{ATP} = 2\text{ produced} - 2\text{ consumed} = \mathbf{0\text{ ATP}}$.

\item \textbf{(A)} Both statements are correct. In fermentation, cytosolic $\text{NADH} + \text{H}^+$ is reoxidized directly by reducing pyruvate (or acetaldehyde). In aerobic respiration, $\text{NADH} + \text{H}^+$ transfers its electrons to the mitochondrial electron transport chain (via shuttles), with $\text{O}_2$ as the final acceptor.

\item \textbf{(A)} Water displaces air from soil pores, depriving roots of $\text{O}_2$. Roots shift to alcoholic fermentation, yielding only $2\text{ ATP}$ per glucose (insufficient to power active mineral uptake and root maintenance) while accumulating toxic acetaldehyde and ethanol.

\item \textbf{(C)} In glycolysis, $2\text{ molecules}$ of $\text{H}_2\text{O}$ are released during the dehydration of two 2-phosphoglycerate molecules by enolase (and $2\text{ H}_2\text{O}$ are released during the formation of $2\text{ ATP}$ from $2\text{ ADP} + 2\text{ P}_i$, while $2\text{ H}_2\text{O}$ were used in the preparatory phase). In the standard biochemical net reaction ($\text{C}_6\text{H}_{12}\text{O}_6 + 2\text{ADP} + 2\text{P}_i \rightarrow 2\text{C}_2\text{H}_5\text{OH} + 2\text{CO}_2 + 2\text{ATP} + 2\text{H}_2\text{O}$), **$2\text{ H}_2\text{O}$** are net produced.

\item \textbf{(B)} *Leuconostoc mesenteroides* carries out heterolactic fermentation (via the phosphoketolase pathway), fermenting glucose into a mixture of lactic acid, ethanol, and $\text{CO}_2$, which initiates sauerkraut and kimchi fermentation.

\item \textbf{(A)} The **Crabtree effect** occurs when high glucose flux overwhelms the respiratory capacity of mitochondria (via catabolite repression of respiratory enzymes and saturation of pyruvate dehydrogenase), diverting excess pyruvate into pyruvate decarboxylase and ethanol production even in the presence of $\text{O}_2$.

\item \textbf{(A)} Pyruvate decarboxylase cleaves a $\text{C}-\text{C}$ bond without hydrolysis or oxidation, classifying it as a **Lyase** (EC 4). Both Alcohol dehydrogenase and Lactate dehydrogenase catalyze redox reactions involving $\text{NAD}^+/\text{NADH}$, classifying them as **Oxidoreductases** (EC 1).

\item \textbf{(B)} Lactic acid dissociates into lactate and $\text{H}^+$ ions. The resulting drop in sarcoplasmic pH (acidosis) inhibits rate-limiting glycolytic enzymes ($\text{PFK-1}$) and reduces the sensitivity of troponin to $\text{Ca}^{2+}$, causing muscle fatigue.

\item \textbf{(B)} Fructose-1,6-bisphosphate bypasses both $\text{ATP}$-consuming steps of the preparatory phase ($0\text{ ATP}$ consumed). Each $\text{F-1,6-BP}$ splits into two trioses, yielding $4\text{ ATP}$, $2\text{ CO}_2$, and $2\text{ Ethanol}$. For $4\text{ molecules}$ of $\text{F-1,6-BP}$: $4 \times 4 = \mathbf{16\text{ ATP}}$, $4 \times 2 = \mathbf{8\text{ CO}_2}$, and $4 \times 2 = \mathbf{8\text{ Ethanol}}$.

\item \textbf{(D)} Statement (D) is incorrect because NCERT explicitly states that $\text{NADH}$ is oxidized to $\text{NAD}^+$ **rather slowly** in fermentation, whereas the reaction is **very vigorous** in aerobic respiration.

\item \textbf{(C)} Aerobic respiration produces $38\text{ ATP}$ per glucose, while fermentation produces $2\text{ ATP}$ per glucose. To maintain the same $\text{ATP}$ supply rate, glucose consumption must increase by $38 / 2 = \mathbf{19\text{ times}}$.

\item \textbf{(A)} *Clostridium acetobutylicum* performs **ABE fermentation**, producing **Acetone, Butanol, and Ethanol** (historically vital for synthesizing cordite during World War I).

\end{enumerate}

\end{document} 
</USER_REQUEST>
<ADDITIONAL_METADATA>
The current local time is: 2026-11-16T11:56:29+05:30.

The user's current state is as follows:
Active Document: /Users/laxmikumari/Desktop/web/testseries/scripts/add_botany_zoology_to_part_11.mjs (LANGUAGE_UNSPECIFIED)
Cursor is on line: 1
Other open documents:
- /Users/laxmikumari/Desktop/web/testseries/.env.local (LANGUAGE_UNSPECIFIED)
- /Users/laxmikumari/Downloads/neet_five_kingdom_classification_test.tex (LANGUAGE_UNSPECIFIED)
- /Users/laxmikumari/Downloads/GOC_Extracted_Questions.tex (LANGUAGE_UNSPECIFIED)
- /Users/laxmikumari/Downloads/neet_biology_question_bank_phases_of_growth_set_2.tex (LANGUAGE_UNSPECIFIED)
- /Users/laxmikumari/Desktop/web/testseries/src/app/api/admin/notifications/route.js (LANGUAGE_UNSPECIFIED)
</ADDITIONAL_METADATA>